\documentclass[12pt]{article}

\usepackage{hyperref}
\usepackage{natbib}
\bibpunct[: ]{(}{)}{;}{a}{}{,}
\usepackage[margin=1in]{geometry}
\usepackage{setspace} 
\usepackage{graphicx}
\usepackage{lscape}
\usepackage{subfigure}

\hypersetup {
	pdfauthor={Omar Lizardo},
     pdfpagemode = UseNone,
     pdfstartview=FitH,
    	citebordercolor = 1 0 0,
	}
	     
\begin{document}

\title{Cultural Capital in Interaction\footnote{Direct correspondence to Omar Lizardo, University of Notre Dame, Department of Sociology, 810 Flanner Hall, Notre Dame, IN, 46556  Tel: (574) 631-1855, Fax: (574) 631-9238, Email: \href{mailto:olizardo@nd.edu}{\tt olizardo@nd.edu}.}}

\author{Omar Lizardo\thanks{I would like to thank Gene Halton, Robert Fishman and Mary Ellen Konieczny for helpful comments on earlier drafts.}
\\\textit{University of Notre Dame}}

\date{\today}
\maketitle

\begin{abstract}
In this paper I analyze survey data from a probability sample of Americans to  for the first time document quantitatively the distribution and predictors of what I call ``culture-mediated interaction'': essentially, the use of cultural content to generate conversation with friends.  I argue that this phenomenon--central to many {\em theories} of the connection between culture and stratification--deserves more {\em empirical} attention. I show that culture-mediated interaction is a pervasive feature of the lives of many Americans, especially those who are culturally active. In addition, I show that rates of culture-mediated interaction appear to be stratified by the usual markers of status (such as educational attainment), and for some cultural domains, this stratification is independent of frequency of cultural engagement.  I close by outlining the theoretical and empirical implications of taking the phenomenon of culture-mediated interaction seriously. 
\end{abstract}
\onehalfspacing

\section{Introduction}

Most of the contemporary research at the intersection of the sociology of culture and stratification is mainly concerned  with exploring inequalities linked to social position in the relative rates of access to and engagement with symbolic goods  \citep{dimaggio00, vaneijck00, katzgerro04}.  The central issue  remains the extent to which taste and culture consumption oriented lifestyles continue to be associated with social position \citep{tomlinson03}.  This common concern can be appreciated whether the construct of ``social position'' is conceptualized in terms of educational attainment \citep{bourdieu_darbel91, bryson96, vaneijck99}, membership in broadly defined occupational groups \citep{bourdieu84, dimaggio_useem78a, dimaggio_useem78b, peterson_simkus93, tomlinson03}, aggregate measures of social class \citep{katzgerro_shavit98, katzgerro02, katzgerro06}, or constructed scalar indices of social status based on social linkages among and within occupational groups \citep{chan_goldthorpe05, chan_goldthorpe07a, chan_goldthorpe07b}. 

The overall conclusion reached in this still vibrant line of research is, that contrary to the predictions of ``postmodernization'' theories that foresaw the decline in the linkage between class and lifestyle in market-dominated ``consumer societies'' lifestyle practices continue to be strongly shaped---albeit in increasingly complex ways \citep{peterson05}---by social position \citep{holt98, katzgerro99, tomlinson03, tampubolon08, warde_etal08}.  Contemporary forms of cultural stratification continue to be structured according to markers associated with social status and education rather than economic class \citep{alderson_etal07,chan_goldthorpe07a, chan_goldthorpe07b, katzgerro02, katzgerro04, sullivan_katzgerro07, silva06, bennett_etal99}.  Confirming the basic findings of an earlier generation of research \citep{baumol_bowen66, dimaggio_useem78a, ganzeboom82}.

The strong focus of the cultural stratification literature on the question of {\em whether} lifestyle practices continue to be unevenly partitioned according to class and status markers, has resulted in a relative lack of attention to an equally important question:  what is the primary {\em role} (or ``function'') of aesthetic consumption in the everyday lives of members of different social strata?  More precisely, given that some groups---usually those who enjoy markers of privilege and status---are more likely to engage culture and the arts more than other groups, to what {\em uses} do they put the culture that the consume?  The relative neglect of the behavioral consequences of aesthetic consumption  appears even more surprising in light of the fact that the issue of the role played by engagement with symbolic goods was an continues to be central to the theoretical corpus from which most analysts of the culture-stratification linkage depart in the first place \citep{collins75, dimaggio87, douglas_isherwood96, fine77}.


A key insight to be gained from this (largely theoretical) tradition is that if we cannot get a solid theoretical handle on the uses to which persons put the cultural repertoires acquired through aesthetic experiences it is difficult to analytically justify their status as ``cultural capital'' in a given social setting \citep{mohr_dimaggio95, dimaggio96, dimaggio00, carter03}.  This is the reason why in their influential paper Lamont and Lareau \citeyearpar{lamont_lareau88} linked the notion of cultural capital to an explicit conceptual framework that highlighted the active deployment of a given repertoire of cultural tastes and lifestyle practices for the purposes of drawing symbolic boundaries. In a similar way, in key theoretical statements DiMaggio \citeyearpar{dimaggio87, dimaggio00} is clear---building on the early work of Collins \citeyearpar{collins75}---in connecting the aesthetic consumption of cultural goods to routine exploitation of cultural knowledge in social interaction within everyday settings (see also \citet{erickson96} and \citet{lizardo06a} for empirical applications of this idea).  It is thus safe to say that the set of cultural aptitudes that result in the relatively higher likelihood to consume cultural goods only becomes actualized as cultural capital when it becomes linked to everyday processes of informal socialization and thus status group formation in the Weberian sense \citep{lizardo06a}.

In this paper, I draw out the empirical implications of this set of starting ideas for our understanding of the {\em social uses} of aesthetic consumption in interaction.  My analytical departing point is that the everyday use of cultural knowledge and experiences gained through aesthetic consumption can be thought of as a social resource that becomes the primary fodder for social interaction in conversation in both relatively formal \citep{long03} and informal \citep{fine77} settings. Inequalities in cultural access can then be understood as reflecting inequalities in the ability and opportunity to engage in certain forms of informal interaction that Simmel \citeyearpar{simmel49} referred to by the name of ``sociability.''  I argue that the primary uses of aesthetic consumption are social \citep{dimaggio87}.  Even if the ``gratifications'' that the consumption of aesthetic goods provides can sometimes be individual, private and  \citep{denora00, hennion01, holbrook_hirschman82}.

This has important implications for our understanding of the link of culture and stratification in late-modern societies.  If the argument offered here is on the right track it becomes clear that for culturally-advantaged members of high-status groups the deployment of cultural capital {\em in interaction} \citep{eliasoph_lichterman03} during face-to-face talk serves as the primary {\em relational mechanism} \citep{tilly98} that links social position to aesthetic consumption.  This mechanism explains both status-based advantages premised on the maintenance of (homo)social bonds, as argued in the classic Weberian tradition, as well as those social advantages tied to the ability to connect across distant social positions in the social structure \citep{erickson96, lizardo06a}.

\section{Culture as a social resource:  Specifying the phenomenon}
The empirical issue that is the focus of this paper pertains to the ``behavioral implications'' of the aesthetic consumption of popular (and fine arts) culture \citep{fine77}.  This phenomenon can best be characterized through the frequently noted observation that cultural knowledge obtained from aesthetic consumption experiences will become the primary content used for the initiation and maintenance of conversational interaction \citep{dimaggio87, collins75}.  I refer to this (presumably) pervasive phenomenon as {\em culture-mediated social interaction}. Culture-mediated interaction can technically defined as the deployment of cultural knowledge obtained from aesthetic consumption experiences as a resource to generate conversation. In a nutshell (and less technically): the discussion of ``cultural matters'' with friends.

Culture-mediated social interaction is a primary example of ``sociability'' in late-modern societies, and a key marker of social status \citep{collins75, lizardo06a}.  Sociability is social interaction that is not driven by explicit instrumental purposes, but that is keyed to the continuing repetition of episodes of social intercourse for their own sake \citep{simmel49}. This is a form of social intercourse explicitly dissociated from direct instrumental pursuits, as a form of association ``which does not have a strictly economic or business purpose'' but access to which (and restrictions on) are a key marker of the status group boundaries \citep[114]{weber94}. It is in this respect that sociability can be thought of as the ``play form'' of the more general process of association (coming together from some purpose).  This interaction cannot be driven by traditional ``instrumental'' contents (e.g. work, politics, etc.), but requires instead ludic content that serves to simply sustain the focus of attention in the interaction itself \citep{collins75}. I argue that the ``entertainment talk'' typical of those who are active culture consumers is one of the forms of cultural knowledge that is most apt for these playful purposes.
        
This Simmelian mechanism is at the center of Collins' microinteractionist theory of culture and stratification.  According to Collins, 
        
\begin{quotation}
[w]ith conversation for enjoyment, we have left the serious realms and are now fully inside the world of sociability.  The social relationships involved also change.  Practical talk establishes relationships for a very limited purpose \ldots alliances based on extrinsic rewards, not on personal liking \ldots Entertainment talk operates on a different level.  It is talk for the sake of talking, as an enjoyment in itself.  Here the talk is embedded in nonverbal rituals that are even more important than the words themselves \ldots People who participate in it to any substantial degree are friends of some sort \ldots \citep[124]{collins75}.
\end{quotation}

Following this line of reasoning we can begin to appreciate with clarity the ``social'' function of symbolic goods consumption. These functions are intimately related to the maintenance and enhancement of status-linked social cohesion in small groups \citep{fine77}.  In this respect, Fine's basic complaint in relation to research of popular culture in the early 1970s, can be repeated vis a vis research on cultural stratification today.  This research, generally ignores ``the behavioral implications of \ldots culture [consumption], particularly as these implications relate to social interaction'' \citeyearpar[453]{fine77}.  Thus, if culture results in talk, and the micro-realities of stratification and status are reproduced through talk, then culture is intimately implicated in the everyday, mundane reproduction that most researchers interested in stratification dynamics study from more Olympian heights \citep{collins04}.  

It is precisely because ``culture produces considerable discussion, and thus \ldots[its] content \ldots will become a resource in conversational interaction'' \citep[453]{fine77} that understanding the phenomenon of culture-mediated interaction should be a primary goal of those interested in the linkage between culture, stratification and segmentation of social interaction by status.  Accordingly, those cultural goods which find an audience, will be deployed as the basis of conversation and will read the ``fodder for least-common denominator talk.'' We should expect that the consumption of aesthetic goods infuses local conversation-based interaction rituals with a common focus of attention allowing persons to spend time engaged in (Simmelian) interaction for its own sake \citep[443]{dimaggio87}. For instance, Elizabeth Long \citeyearpar{long03} in her study of women's reading groups in Houston, Texas, finds that cultural knowledge and taste become the primary content of social interaction. Consistent with the claim that connects the consumption of aesthetic products and sociability, she finds that reading group members ``tend to press books into service for the meanings that they transmit and the conversations they generate \citeyearpar[73]{long03}.  

It is safe to conclude, in agreement with an earlier tradition on the behavioral consequences of media and the arts, that the consumption of cultural goods has many ``effects.'' But these effects go beyond its perceptual, psychological and emotional impact at the point of direct reception.  The main effect of culture is social, in terms of producing content that becomes ``processed'' in face-to-face interaction.  We can then think of ``[t]he interacting group \ldots as a culture processing machine, which takes popular culture as experienced by group members, and transforms it into content useful for the continuation of interaction'' \citep[460]{fine77}.  I propose that in the contemporary context, aesthetic experiences serve the role of providing fodder for ``entertainment talk'' \citep{collins75, dimaggio87, fiske87, frith98}, and that the primary uses to which individual put the knowledge gained through experiences of aesthetic consumption are  ``social'' in the aforementioned sense.  

Thus characterized, the key empirical dynamic that is the focus of this paper concerns the phenomenon of conversation that involves references and allusions to arts and culture as its primary content.  We should expect that arts consumption will be associated with a higher frequency of these types of culture-mediated social interaction. In addition we should also expect that the relative frequency of these opportunities for culture-mediated interaction will be stratified along the same lines that culture consumption and social connectivity are stratified \citep{dimaggio87, lizardo06a}.  This means that individuals who are already have a higher predisposition to engage in aesthetic consumption---indexed by levels of schooling \citep{bourdieu84, bourdieu_darbel91}---will also be more likely to themselves utilize the cultural knowledge gained from this experience to engage in social interaction.

\section{Theoretical and Substantive Importance of the Phenomenon}       
 Why should the phenomenon of culture mediated interaction be of interest to sociologists interested in the linkage between culture and stratification?  While the connection between aesthetic consumption and social interaction (and thus the uses of culture to build and maintain social networks) has been a relatively neglected issue in the current sociology of culture literature the phenomenon has been obviously perceived to be a key issue in some of the foundational theoretical statements in the field. From this literature we can garner a series of leads as to why we should find that culture-mediated interaction is relatively generalized and why it should be stratified by the familiar markers of status.  
                   
\subsection{Relationship formation}
Culture-mediated interaction allows for the formation and renewal of social relationships.  This is a ``need'' that is structurally induced by the fact that, as recent research on the dynamics of personal networks show, new relationships are constantly being formed and old ones renewed (or dropped) during the adult life-cycles.  The most surprising thing about ``intimacy'' and ``friendship'' is thus, how unstable it tends to be \citep{wellman_etal97}.  Cultural knowledge gained through aesthetic consumption is a key resource with which to form {\em new} social relationships.  ``Culture talk'' is especially useful at the \emph{initial} stages of ``the acquaintance process'' \citep{newcomb61}. Culture-mediated interaction is helpful in this respect because it allows us to search for ``deeper'' areas of commonality and agreement without much risk (as would inquiring about politics or economic position), so that the interaction can proceed smoothly.  Culture mediated interaction is a lubricant that keeps conversation going when intimate content is not yet available.   When individuals do not have a previous social or cultural bases for interaction based on common membership in the same social network or common biographical experiences they instead rely ``\ldots upon common cultural topics on which to build a conversation.''   Cultural goods and aesthetic experiences in general  provide ``\ldots a topic which is both interesting and non-controversial.  It provides an entr�e to interaction'' \citep[460]{fine77}. 

Culture-mediated interaction that makes use of cultural knowledge about symbolic goods that are truly popular may cut across class and status lines may in this way serve to connect persons across different positions in the social structure and across different taste publics, facilitating coordination across hierarchical positions \citep{erickson96}. It is precisely because ``[t]he content and symbolic meaning of these products has diffused sufficiently widely as to make them a viable conversation topic for strangers and new acquaintances'' \citep[460-461]{fine77}, that it  allows persons to create and sustain social connections--even if sometimes brief and sporadic (so-called ``weak ties'')--when there exist large socio-demographic differences between them \citep{lizardo06a}.

\subsection{Status-based solidarity in small groups}
Culture-mediated interaction serves as the primary mechanism to generate and reproduce the micro-solidarities that hold status-groups together \citep{ostrower98}.  The common knowledge gained through culture talk serves as a cognitive and emotional basis of solidarity within the micro-ritual generated in the interaction, as persons become entranced by ``the joy of sharing names'' \citep{dimaggio87}. Culture-talk``\ldots facilitates the sociable intercourse necessary for acquaintanceships to ripen into friendships'' \citep{dimaggio87}.   

Status-based exclusion occurs largely through micro-mechanisms governing entry (or determining systematic exclusion from) class-structured micro-worlds organized around specific conversational interactions \citep{collins75, collins00}.  The experiential and phenomenological ``stuff'' of stratification is essentially based on talk, and much of this talk is about culture.  This is the reason why we can consider taste to be ``\ldots a form of ritual identification and a means of construction social relations\ldots It helps to establish networks of trusting relations that facilitate group mobilization and the attainment of such social rewards as desirable spouses and prestigious jobs''  (DiMaggio 1987:  443).  Culture-mediated interaction thus serves to increase the cohesion and social closure of small groups, providing a way to segment them from the larger social structure;  this is why we can think of culture as being much more ``\ldots than just an acquaintanceship device.  Member know each others' backgrounds, and have a general idea of others positions, values, and beliefs.  In this case popular culture can serve as a mechanism for establishing, confirming and reifying group beliefs and attitudes'' \citeyearpar[461]{fine77} 
                   
The primary function of aesthetic consumption experiences is precisely to serve as the non-instrumental content that justifies the existence of status-segregated conversational rituals:  ``\ldots talk about entertainment-the discussions of movies, music, sports, actors and athletes\ldots make up so much of casual chitchat.'' This becomes even more important ``\ldots if we include talk about things people consume or seek out-food, clothes, houses and gardens, restaurants scenery, vacations.  In some social circles, this entertainment talk makes up most of the content of conversations; in all classes, it is important" \citep[125]{collins75}.  The Simmelian language of ``social circles'' is not accidental, since as Kadushin \citeyearpar[797]{kadushin66} noted in his classic study on the subject, Simmelian urban social cirles are produced and reproduced through shared {\em cultural tastes}--e.g. concert going, visiting the museum, etc.--and other activities intrinsically connected to culture-mediated interaction.                     
                   
\subsection{Linkages within loosely bounded and sparsely knit networks}     
Culture-mediated interaction is not only  the ``bond'' that keeps status groups together, but also the ``solvent'' that allows for sporadic connections across status lines and that provides new forms of high-status group solidarity not built around traditional notions of commonality anchored to geography, community and shared institutional affiliations and kinship structures \citep{holt98}.  One of the most significant transformations in the social logic that governs status in contemporary (post)industrial societies consists of the rising role of the arts, and mass produced culture in general, in providing the ``baseline'' forms of cultural knowledge necessary to sustain interaction for persons who cannot come to expect commonalities based on traditional symbols of acquaintance and recognition. 

This shift has occurred in the context of the social transformation of status-based collectivities from solidary groups with a strong allegiance to community and place, to looser, wide-spanning ``personal networks'' characterized by increasing geographic and social mobility \citep{wellman_etal88}.  This process acquires more importance as these network ties have been transformed in the contemporary context of increasing geographic mobility \citep{griswold_wright04} and the decline of the traditional bonds characteristic of traditional local communities \citep{wellman79}, which used to be centered around kinship and spatial contiguity \citep{dimaggio_mohr85}.  As the personal networks of socially mobile members of the upper-middle class have become more fluid, loosely knit, more discontinuous and less tied to spatial contiguity and kinship \citep{wellman_wortley90, wellman_etal88}, the arts and popular culture have come to increasingly serve as the ``default" forms of knowledge that connect people across different interactin foci.
                   
Social interaction based on culture comes to occupy center stage because aesthetic consumption experiences are an ideal ``content'' to fuel the process of Simmelian sociability. In contrast to material goods which ``are physically present and visible'' and which are more likely to produced inviduous distinction, cultural consumption, ``\ldots is invisible once it has occurred.  This evanescent quality makes artistic experience, described and exploited in conversation, a portable and thus potent medium of interactional exchange.'' \citep[443-444]{dimaggio87}.  Popular culture and the arts (along with ``fashion, cuisine and sport'' are the only analog to a common cultural currency in complex societies \citep[443]{dimaggio87}.  The arts are not only a generator of cultural knowledge and are also one of the most important facilitators of culture-mediated interaction.                 
            
In sum, some of the most important lines of theory in the culture and stratification literature point to the central importance of the phenomenon that we have labeled and characterize as {\em culture mediated social interaction}.  These theoretical streams point to this phenomenon as a pervasive, ubiquitous feature of late modern societies.  According to this line of theory, culture-mediated interaction is key to understanding the dynamic, interactional origins of the social discontinuities and status-based social boundaries that characterize the modern social order \citep{illouz97}. This makes its almost total neglect in the more empirically oriented recent literature in culture and stratification an even more puzzling oversight. 
                   
The formal and informal system of production of aesthetic goods provide as ``input'' the very content that serves as the scaffolding around which everyday episodes of ``talk'' and sociability are built.  The socially segmented micro-orders based on face-to-face (and increasingly electronically mediated) interaction require a constant stream of new interactional content in order to be sustained and thus reproduce themselves over time. It is thus no surprise, as we saw above, that the stratification of culture-mediated interaction should mirror in non-trivial respects patterns of stratification found in the consumption of aesthetic goods and experiences. This is consistent with Simon Frith's \citeyearpar[4]{frith98} suggestion that ``[p]art of the pleasure of popular culture is talking about; part of its meaning is talk �''   The micro-solidarities that form the basis of status-linked orders but also form of the core of what has come to be called ``social capital'' are may thus be built around the most seemingly innocuous forms of conversational subject matter.  This implies that the phenomenon of culture-mediated interaction is relevant for those interested in the connection between cultural and social capital \citep{lizardo06a, jeannotte03, dimaggio96}.  As John Fiske's observed, civic membership and ``belonging'' to neighborhoods and communities is based on talk, and ``[m]uch of this talk is about the mass media and its cultural commodities'' \citeyearpar[73]{fiske87}.
                
\section{Empirical Analysis}
\subsection{Data Source}
In what follows I examine data taken from the Arts and Culture Supplement to the National Social Survey fielded in 2002 and 2004.   The NSS is a weighted (by age, education, marital status, gender, race, ethnicity) probability sample respondents obtained through computer-aided telephone interviews. This survey is particularly appropriate for our purposes because it contains information not only about culture consumption experiences, but also about the likelihood that individuals engage in informal social interaction featuring cultural content taken from those aesthetic consumption episodes.  Because the items of interest were fielded in both waves, I pool the samples for which I have complete data on all socio-demographic variables from the 2002 ($N=955$) and 2004 ($N=776$), which leaves us with a total combined sample of 1,731 respondents. 

\subsection{Culture-Mediated Social Interaction Survey Items}

As part of the survey, respondents were asked the following question:

During the past twelve months, approximately how often have you had a discussion with someone you know about any of the following:
\begin{enumerate}
\item  {\em A book, a poem, or a story?}
\item  {\em A musical performance, live or recorded.}
\item  {\em A movie?}
\item  {\em A play or other dramatic performance?}
\item  {\em An art show or work of art?}
\end{enumerate}

Respondents were asked to choose from the following response categories: 1) very often, 2) fairly often, 3) not very often, 4) or not at all. 
                
\subsection{Aesthetic Consumption Survey Items}
                                            
Respondents were also asked to report the frequency with which they engage in various forms of aesthetic consumption:
 
\begin{enumerate}
\item  {\em Books}:  Approximately how often during the past 12 months, have you read a book?
\item  {\em Movies}:  Approximately how often during the past 12 months, have you been to the movies? 
\item   {\em Music}: Approximately how often during the past twelve months would you say you listened to music at home, for instance on CDs, tapes, records, a computer, or on the radio?
\item  {\em Live performances}: Approximately how often did you go to                 any live performances such as plays, concerts, musicals, or dance?
\item  {\em Art}:  Approximately how often did you go to any art shows or museums during the past 12 months? 
\end{enumerate}  

\section{Analysis and Results}
\subsection{Distribution of culture-mediated interaction in the United States}
\subsubsection{Overall frequency}
I begin by attempting to provide an answer to a simple question: what is the relative distribution of the frequency of culture-mediated interaction in the contemporary United States? Insofar as culture-mediated interaction is a key process in current theories of the culture/stratification link, we should expect it to a fairly widespread phenomenon.  Because there has been virtually no quantitative research on issue however, it is hard to estimate with certainty how much is "widespread."  Given what we know from the few qualitative studies that have been conducted, we can  expect to observe close to a majority of the population reporting high or moderately high frequencies of culture-mediated interaction.\footnote{As with any study that uses arts participation data in order to draw larger conclusions about the operation of culture in society at large, the big caveat is of course that any observed frequency of culture-mediate interaction is probably an underestimate, as persons can only report interacting with friends on the basis of the cultural forms included in the survey.  Clearly, persons who are observed to not interact with friends very often on the basis of content taken from the cultural domains included in the survey, may still be doing so on the basis of cultural forms not included in the survey.}
      
Figure~\ref{fig1a} shows the univariate distribution of culture-mediated interaction across the five broad realms of cultural activity included in the 2002 and 2004 National Social Surveys.\footnote{the circle is the observed percentage and the vertical line represents the 95\% confidence interval, calculated using the ``exact'' binomial estimate for the observed proportion.}  The distribution of culture-mediated interaction frequencies across all five cultural domains supports the expectation that persons use culture as a discussion topic in a fairly regular manner.  78\% of respondents report frequently (either ``fairly'' or ``very'' often) using {\em at least one} of the cultural domains included in the survey as a subject matter in conversation with friends.
                
When it comes to the three most popular cultural domains  (movies, literature and music), it is clear that a majority (or close to a majority) of respondents report having discussed matters related to that cultural domain with a friend.  About 60\% of persons report talking to friends about movies ``very often'' or ``fairly often''; the corresponding percentages for books and music are 52\% and 46\%. Even for the two cultural domains which are expected to have the most exclusive audiences (in terms of such socio-demographic characteristics as education and age), in this case theater and the arts, 26\% and 19\% of respondents report relative frequent rates of culture-mediated interaction respectively. 
                
\subsubsection{Frequency among cultural active respondents}
The raw estimates of culture-mediated interaction give a somewhat misleading picture of the extent to which the phenomenon occurs. It is reasonable to suppose that a necessary precondition for persons to deploy any type of cultural knowledge during interaction is that they have acquired this knowledge.  It is a good bet that knowledge about symbolic goods is acquired through direct engagement with the formal and informal culture industries (and non-profit organizational fields) that deliver these goods.  Thus, we should expect the relative chances of observing high rates of culture-mediated interaction to be mediated by the extent to which the persons engages the cultural field in question as a consumer.
      
Figure~\ref{fig1b} shows the conditional univariate distribution of culture-mediated interaction among those respondents who report engaging in the corresponding cultural activity either ``fairly often'' or ``very often.''  I find that not surprisingly, rates of culture-mediated interaction among active consumers of the cultural forms in question are substantially higher than the average.  About 90\% movie, 71\% of literature and 51\% of music consumers report engaging in relative frequent rates of culture mediated interaction.  The most dramatic relative increases in relation to the average is observed for the two cultural domains associated with the non-profit ``fine arts'' field:  53\% of active dramatic theater consumers and 61\% of active art consumers report discussing this type of cultural content with a friend.  
      
Interestingly, the univariate distribution for culture-mediated interaction based on music shifts the least after consumption is taken into account.  This indicates that this type of cultural content is deployed routinely without necessarily being tightly connected to consumption.  It can also mean that the ``uses'' of music as cultural form may be unique:  music may constitute both a ``social'' resource for the cementing of connectivity in social relationships, and also a private ``resource'' for personal enjoyment and the customization of ``life-spaces'' \citep{denora00}.
      
Another noticeable fact is that both movies and literature seem to be ``inherently social'' cultural domains in these data.  It thus appears to be almost (statistically at least) ``non-sensical'' to be an active movie consumer or book reader without using this cultural knowledge in interaction.  The concrete mechanisms governing the way in which these two cultural forms are deployed in everyday settings may be distinct however. The fact that in these data the proportion of respondents who actively engage the cinematic field {\em without} at least occasionally using this cultural knowledge in conversation corresponds well with Fine's \citeyearpar{fine77} original proposal that cultural goods with widespread appeal would also be the most likely to be mobilized in interaction in small groups. Research on book groups by Long \citeyearpar{long03} and Illouz's \citeyearpar{illouz03} research on Oprah's televisual ``book club'' suggest that books are a primary conversational resource among a wide variety of audiences.  However, literature conversations tend to occur within the relatively formalized ambit of the formal groups purposefully designed for the purposes of literary discussion (for a history of the book club institution in the U.S. see Long \citeyear{long03}).  This is in keeping with the highly institutionalized perception of (most) literature (and reading as a cultural practice) as ``serious.''   Conversations about such cultural fare as movies, however, serves to cement interaction in less formalized social settings \citep{fine77}.
      
\subsection{The stratification of culture-mediated social interaction}
While most persons are indeed likely to engage in some form culture-mediated interaction, it is to be expected that as a general rule, culture-mediated interaction should be stratified along indicators of status and prestige.  We should expect this to be the case for two reasons:  first, culture consumption is itself stratified, with more educated individuals being more likely to engage in a wide range of aesthetic consumption activities than the less educated \citep{bourdieu84, bourdieu_darbel91}, and as we have already seen, culture-mediated interaction is much more likely to be observed among active culture consumers than among inactive ones.  
                
Second, opportunities for social interaction and such indicators of social connectivity as network size and proportion non-kin in the typical discussion networks are stratified by the same types of status and prestige indicators, in particular educational attainment \citep{marsden87, fischer82, burt90}.  Insofar as culture mediated interaction is more likely to take place among members of informal cliques, we should expect that those persons most likely to have wide-spanning, loosely knit networks of acquaintance, composed of elective relationships of ``friendship'' (in contrast to ascriptive relations of kinship) should be more likely to engage in culture-mediated interaction.  Most research on personal discussion networks indicates that the higher educated are more likely to enjoy these types of local relational environments than the least educated \citep{burt90, marsden87, erickson96, lizardo06a, fischer82}.      
                
Figure~\ref{fig2a} shows the predicted probabilities by education levels from a logistic regression predicting whether the respondent reports discussing movies, books, plays, the arts or music with a friend during the past year. The regression model is estimated using a specification that holds constant respondent's age, gender, race, marital status, urban location and year of survey.  The results are fairly unambiguous: across all fives cultural domains, as education increases, so does the odds of having had a discussion with a friend involving culture. This means that, as expected, culture-mediated interaction is a stratified phenomenon:  the educated are more likely to deploy this cultural knowledge in interaction. 
     
How does variation across levels of education in regards to culture-mediated interaction compare to variation on the same outcome across the levels of other indicators of social position?  The top panel of Table~\ref{tab1} shows how variation across levels of education compares to heterogeneity across the levels of the two other variables included in the model which proved to have the strongest statistical contributions on the models predicting culture-mediated interaction:  age (in seven categories) and gender (in two categories, male and female).  For each equation I take the model $\chi^{2}$, which is an indicator of the total explanatory power supplied by the systematic part of the model, and calculate what proportion of it is due to each set of variables (e.g. a set of four dummy-coded variables for education with less than high-school as the reference, a set of six dummy-coded variables for age group with the 18-24 age-group as the reference and a dummy-coded variable for gender with males as the reference).  The first three columns show the $\chi^{2}$ test for the hypothesis that the coefficient corresponding to the four education dummy variables is equal to zero.  And the same is shown in the subsequent six columns for age and gender.  The last three columns show the proportion of the total model $\chi^{2}$  for which each variable is responsible.  The results show that except for culture-mediated interaction involving movies (for which age is by and large the main explanatory factor), education accounts for a relatively substantial proportion of the explanatory power when it comes to literature, music, theatrical plays and the arts.  Among these last two genres, the large majority of the model's explanatory power is due to education. 

\subsection{Variation in aesthetic consumption across social position}
How does variation across levels of education, age and gender on the propensity to engage in culture-mediated interaction compare to variation across the same three social position indicators on the frequency of aesthetic consumption in the same five cultural domains?  The bottom panel of Table~\ref{tab1} shows the results.  I find that for the most part, the results mirror those obtained when frequency of culture-mediated interaction (rather than engagement) in that domain is the outcome.  When it comes to relative rates of engagement in the higher-status cultural pursuits (literature, the arts, drama), education accounts for the bulk of the explanatory power of the model (57\%, 84\% and 66\% respectively), while for movies age ($\Chi^{2}=175.7, p<0.01$) is the primary factor differentiating avid consumers from abstainers.  In addition, variationa accross gender only figures prominently in the case of reading \citep{tepper00, long03}.

There is one key exception to this pattern of results. In the case of the frequency of culture-mediated \emph{interaction} music mirrors the relatively high-status domains in displaying more variation across levels of education than it does across age-levels of gender.  However, when it comes to frequency of music \emph{consumption} music mirrors the more popular domain of movies in displaying no statistically discernible variation across education levels ($\Chi^{2}=4.3, p=0.37$), and in displaying the bulk of the variation across age levels instead  ($\Chi^{2}=77.1, p<0.01$). This dissociation between the frequency of aesthetic consumption and interaction involving music is substantively interesting, since it implies that the extent to which music may serve as a conversational resource may vary across musical genres, with high-status genres being more likely to figure as a subject matter in conversation among highly educated strata. In addition, since age is the primary indicator of social position in regards to frequency of engagement in music (averaging across genres), it is possible that among youth, a qualitatively different process of culture-mediated interaction (involving a distinct array of genres) may also obtain.

\section{Social position, aesthetic consumption and culture mediated interaction: Three models of the process}
A key issue that emerges in analyzing the relationship between aesthetic consumption, social position and culture-mediated interaction is the question of their inter-relationship.  This issue is important, because different theoretical representations of the underlying process that generates the association between position, disposition and propensity to engage in culture-mediated interaction imply analytically distinct models of how cultural tastes are shaped by (or shapes) interaction opportunities \citep{lizardo06a, erickson96}.  Contemporary research and theory suggests three (parsimonious) ways of thinking about how they are connected to one another.  Essentially this boils down to the question of establishing which of factor is \emph{temporally and causally prior} to the other ones, and which ones are more likely to function as \emph{intervening mechanisms}. 

While it is a commonplace to acknowledge that causal priority cannot be unambiguously established with cross-sectional data, it is mistaken to conclude that the analyst is therefore helpless. Instead, because all hypotheses that suggest some intuitive arrangement of multiple factors put explicit constraints on a given system of inter-relationships, it is possible to at least ascertain whether the cross-sectional data are compatible to what would be expected by a particular model of the process at ``equilibrium.'' In particular, it is possible to adapt Kendall and Lazarsfeld's \citeyearpar{kendall_lazarsfeld50} influential analysis of spurious correlation for the goal of determining which model is most consistent with a given underlying causal structure \citep{simon54}.   In this section I briefly summarize three views, and propose a simple log-linear models \citep{lever79} that properly characterize them.  

\subsection{Spurious association due to social position}
Under this model, there is no inherent association between aesthetic consumption and culture-mediated interaction.  Instead, the reason why we find that persons who engage the arts are also more likely to talk about the arts, is that those persons can be placed within certain socio-demographic slots which are characterized by recruiting persons why different tastes and different (exogenous, that is unrelated to tastes) opportunities to encounter others and engage in culture-mediated talk.  For instance if the highly educated are more likely to both engage in aesthetic consumption and are also more likely to frequent the types of informal interaction settings in which culture is likely to be brought up as a subject in conversation, then the association between aesthetic consumption and culture-mediated interaction should be greatly reduced \emph{within} levels of education.  

If we imagine a three-way cross-classification between education (with categories $i=1, 2, ...I$), levels of aesthetic consumption (with categories $j=1, 2, ...J$) and frequency of culture-mediated interaction (with categories $k=1, 2, ...K$), this model can be written as:

\begin{equation}
\log(m_{ijk})=\lambda+\lambda_{i}+\lambda_{j}+\lambda_{k}+\lambda_{ij}+\lambda_{ik}
\label{eq1}
\end{equation}

Where $m_{ijk}$ is the expected frequency in the $ijk^{th}$ cell of the table.  In substantive terms, this model says that there is no association between aesthetic-consumption and culture-mediated interaction ($\lambda_{jk}=0|\lambda_{ij},\lambda_{ik}$) within levels of social position (e.g. education or age), which is equivalent to saying that social position accounts for the interaction-taste association.

\subsection{Opportunities for interaction as an intervening mechanism}
An alternative model of the process conceptualizes \emph{interaction opportunities} in certain delimited social foci \citep{feld81} as the primary (intervening) mechanism connecting social position and aesthetic consumption.  This model is more consistent with a network-theoretic view which sees the connection between social position and cultural tastes as a spurious by-product of differentially distributed (across categorically defined social positions) opportunities for interaction in certain settings.  Following Lever's \citeyearpar{lever79} formalization of Blau's  \citeyearpar{blau55} strategy for determining the most likely cross-sectional association structure consistent with this intervening role of interaction opportunities, we can write this model as:

\begin{equation}
\log(m_{ijk})=\lambda+\lambda_{i}+\lambda_{j}+\lambda_{k}+\lambda_{ik}+\lambda_{jk}
\label{eq2}
\end{equation}

Substantively, this model says that the observed association between   social position and aesthetic consumption should be greatly reduced within levels of social interaction ($\lambda_{ij}=0|\lambda_{ik},\lambda_{jk}$).  Thus, those persons who, due to the possession of more expansive and diverse ``personal communities'' embedded in networks or access to social foci in which culture-talk is more likely to occur, are expected to be more likely to acquire and sustain taste for the arts and are also the persons most likely to possess certain markers of status and position in the social structure.

\subsection{Aesthetic consumption as an intervening mechanism}
A third way to conceptualize the process is to construe \emph{cultural tastes} as the primary mechanism joining social position and aesthetic consumption \citep{lizardo06a}.  This conceptualization differs from the network-theoretic model suggested by equation~\ref{eq2} because here tastes are seen as a much more durable part of the person---``embodied'' cultural capital in Bourdieu's \citeyearpar{bourdieu86} terms---and it is through the differential distribution (across social positions) of dispositions towards culture and the arts that we can account for the association between social position and opportunities to engage in culture-mediated interaction. This model implies the following pattern of association:

\begin{equation}
\log(m_{ijk})=\lambda+\lambda_{i}+\lambda_{j}+\lambda_{k}+\lambda_{ij}+\lambda_{jk}
\label{eq3}
\end{equation}

Substantively, this model says that there is no association between social position and the propensity to engage in culture-mediated interaction within levels of aesthetic consumption ($\lambda_{ik}=0|\lambda_{ij},\lambda_{jk}$).  This model is  consistent with a ``culture conversion'' \citep{lizardo06a} argument regarding the process that links social position with culture-mediated interaction opportunities.  Under this model, cultural tastes serve as ``bias'' in contact selection and contact-maintenance \citep{dimaggio93b, vaisey_lizardo}, such that persons with a given set of tastes, tend come to occupy certain social positions (e.g. the rank of the highly educated) and come to surround themselves with similar-taste contacts. Cultural tastes are thus the intervening mechanism connecting a given social position with a higher likelihood of engaging in culture-talk.

\subsection{Results}
Table~\ref{tab2} shows the fit statistics corresponding to the models shown in equations~\ref{eq1}-~\ref{eq3} for the five cultural pursuits included in the NSS. For each cultural domain, I use the social position indicator that is the best predictor of culture-mediated interaction as given in Table~\ref{tab1}. This means that for Books, I rely on both gender and education, and for music I use age and education. This resulted in the estimation of a total of $3 \times (5+2)=21$ log-linear models.  I select the model with the lowest (most negative) BIC score as providing the most likely representation of the data-generating process.

The first thing to notice is that the null model of spurious association by social position is rejected for all cultural domains and for all the indicators of social position. Thus, the association between social position (e.g. education, gender or age) and culture-mediated interaction is not an epiphenomenon produced by the prior association between the propensity to consume culture and the propensity to engage in culture-mediated talk and the person's socio-demographic characteristics.  This also means that, depending on the cultural domain either interaction opportunities can be thought as a mediator---e.g. the ``intervening variable'' \citep{lever79}---between position and cultural dispositions or cultural dispositions are converted (through their linkage to social positions) into opportunities to engage in culture-mediated talk.

The pattern of results shows that depending on the cultural domain and the indicator of social position, the network model or the culture conversion model provide better representations of the data.  The only exception is movies, for which there appear to be no statistical grounds to select between the models, although either the network or culture conversion model offer a large improvement in fit over the null social position model. The network does a better job of accounting for the association between social position, culture-mediated talk and culture consumption when it comes to Books (regardless of whether the social position indicator is gender or education and Music (but only in the case of education).  

The culture conversion model appears to provide a better handle on the association structure between these three variables when it comes to the arts and drama (where the indicator of social position is education) and music (where the indicator of social position is age).  The fact that the culture-conversion model performs best in the case of the highest-status cultural domains is consistent with the results reported by Lizardo \citeyearpar{lizardo06a} who argues (and finds) that high-status culture is most effective in generating larger network of strong ties, and thus is more likely to be implicated in processes of culture conversion generating bounded communities of status-proximate individuals within higher-educated strata.

I believe that there is a theoretically illuminating way to account for these results.  It relies on seeing the case of books as paradigmatic of the operation of an ``embedded social foci'' mechanism accounting for the connection between taste, social position and culture-talk.  As noted above, books are distinctive in being the only cultural domain that has developed a widespread organizational structure in which lay consumers can organize themselves to engage in culture-mediated interaction:  the book club \citep{long03, illouz03}.  While there exist organizations dedicated to the discussion of culture emanating from all of the cultural domains under consideration here, none has achieved the widespread institutionalization among lay audiences as the book club, having remained the purview of relatively selective (and relatively small) circles of enthusiasts \citep{kadushin66}.  Thus, it is not surprising that the ``social foci'' network model outperforms the culture conversion model in the literary domain.  

However, the key piece of evidence substantiating the validity of this interpretation in Table~\ref{tab2} is in the musical domain.  Here, the relative validity of the network model over the culture conversion model is \emph{contingent} on the indicator of social position chosen.  When education is the social position indicator, the network model outperforms the culture conversion model; however, when age is the social position indicator chosen the culture conversion model provides a better representation of the data than the network model.  I submit that the reason why this is the case is that in the case of education, culture-talk involving music approximates the modal case that obtains for books, and for which the network model is the best explanation.  Here discussion takes place within the context of relatively formally organized contexts (social foci) which serve as the primary selection mechanism connecting persons in a certain position with other with similar dispositions.  In the case of age, the results are better interpreted as one in which the social context rather than canalizing tastes towards certain regions of social space \citep{mark98b}, is actually shaped by the distribution of tastes \citep{lizardo06a}, such that contact selection is driven by taste compatibility, and not the reverse.  

This is supported by the results shown in Figure~\ref{fig3}, which shows the predicted probabilities\footnote{The predicted probabilities are obtained from an ordered logit model with frequency of music talk as the dependent variable and dummy-coded indicators of education and age as predictors.} of engaging in culture-talk featuring music as the subject of conversation ``very often'' as they are predicted by Education and Age.  We can see very clearly that (not surprisingly) when it comes to education, culture-talk involving music is the purview of the highly educated; when it comes to age, however, culture-talk involving music primarily occurs among younger respondents. There is no two-way interaction between education and age when it comes to predicting music-based culture-talk in these data ($p=0.29$ for the likelihood ratio test comparing the additive to the unconstrained multiplicative model), which means that education increases music-mediated interaction across \emph{all levels} age. There is an extensive literature that convincingly shows that there exist large differences in the musical tastes of youth and the middle-aged respondents even within levels of education \citep{mark98b, smith94}, so it is a safe bet that highly educated middle-aged respondents are not using the same cultural resources to engage in music-mediated interaction as the young respondents.  There also exists a large literature in cultural studies \citep[e.g.][]{hebdige79, frith98, muggleton00} showing that music-culture plays a key role in constituting and structuring the social life of youth. These studies suggest that processes of conversion of cultural tastes into bounded personal communities (and the use of tastes as a selection mechanism to determine relationships) is more likely to occur under these circumstances.  Among the highly educated, culture-mediated interaction involving music is more likely to similar (but less widespread) to the case of books: occurring in relatively formalized settings that bring together a set of ``enthusiasts,'' who serve as key mediating (and tastek-making) audiences in cultural fields \citep{hannan_polos_carroll07, weber77}.
        
\subsection{The relationship between aesthetic consumption and culture-mediated interaction}
It is to be expected that a key part of the effect of education on culture-mediated interaction is not direct, but it is mediated through the effect of education on aesthetic consumption.  We know that as education increases, so does the likelihood of engaging in aesthetic consumption of all types \citep{warde_etal08}. From Figure~\ref{fig1b}, it can be surmised the culture-mediated interaction is much more likely to be observed among those who are culturally active.
     
However, it is unlikely that education has no direct effect on culture-mediated interaction net of consumption.  This is for two reasons: one theoretical, the other empirical.  First, cultural capital theory leads us to expect that high status individuals are able to demonstrate ``passing knowledge'' of a wide variety of cultural forms \citep{peterson92, peterson_kern96}.  Therefore, their ability to deploy cultural knowledge in interaction should not be completely due to their higher aptitude to consume culture in the first place, but should simply be due to the higher aptitude of the higher educated to be able to produce culture-based talk \citep{collins75, bourdieu84}. 
                
In addition, within broad cultural domains as here defined, we should expect the higher educated to be attracted to those sub-genres the aesthetic qualities of which most lend themselves to verbal expression (mostly due to the presence of ``intellectualized'' critical vocabularies with which to talk about the culture).\footnote{We will explore this hypothesis more fully below.}  Thus, there should be a residual direct effect of education net of consumption, related to the previously demonstrated association between education and the ability to mobilize cultural knowledge in interaction, which itself may be generated by what Bourdieu \citeyearpar{bourdieu84} referred to as the uneven distribution of ``linguistic capital'' across educational groups.
     
The second (more empirical) reason why there should be direct effects of education for at least some class of cultural pursuits is that for some cultural domains, {\em there is no education-based stratification of frequency of cultural engagement that corresponds to the education-based stratification of culture-mediated interaction} observed in figure~\ref{fig2a}.
     
Figure~\ref{fig2b} shows exactly what I mean. These are the predicted probabilities across five educational groups on the frequency of cultural {\em engagement} for the five cultural domains under consideration.  The dependent variable equals one if the respondent reports consuming that cultural forms either ``very often'' or ``fairly often.''  As the figure indicates, while there is a clear education-linked stratification of consumption chances for literature, theater and the arts, there is no education-based stratification in music or movies consumption. Thus, while it is possible that the relatively higher chances of the educated to use cultural knowledge obtained from reading or for patronizing theater and arts disseminating institutions is due to their higher probability of engaging these culture-production fields in the first place, the higher likelihood of the educated to also be able to refer to cultural knowledge related to music and the movies cannot be explained the same factor.
                
Figure~\ref{fig3} provides a more rigorous examination of this claim. The figure shows the ``marginal effect'' \citep{long_freese06}, of education on the probability of frequent culture-mediated interaction from a logistic regression that holds constant age, gender, marital status, urban status and race.  The marginal effect is the predicted change in the probability of the outcome that results from a one-unit change in the dependent variable (in this case it corresponds to making a given educational transition) holding all the other variables (including the chances of having made previous educational transitions) constant.  The ``reduced'' model is a model that does not include the frequency of consumption in that cultural domain as a control, and the full model is one that does. The reduction in the marginal effect of education gives us an estimate of how much of the education effect on culture-mediated interaction using cultural knowledge from a given domain is in fact indirect.
                
The results show that--as expected--frequency of listening to music has a minimal effect in helping us explain the reason why music-mediated interaction is more likely to occur among the educated. The same cannot be said for movies, as holding constant education help us explain about 49\% of the college transition effect and about 37\% percent of the graduate degree transition effect.  That said, it is important to note that even if reduced, education continues to have a significant effect on culture-mediated interaction {\em after} controlling for the frequency with which a person goes to the movies.  
                
The hypothesis that the effect of education on culture-mediated interaction is primarily indirect (e.g. mediated through consumption) fares better when it comes to the most educationally exclusive activities:  reading, patronizing the arts and attending dramatic performances.  In terms of literary consumption, holding constant the frequency with which a person reports reading books helps us account for about 41\% and 47\% of the college and post-graduate degree effects.  When it comes to the arts and theater the effects are even more dramatic, holding constant participation, explains most of the effect of education: in the ``full'' model the $\chi^{2}$ test of the null hypothesis that the effects of education are zero cannot be rejected ($p>0.10$).


                
\subsection{Variation in culture-mediated interaction across genres}
Some aesthetic domains may exhibit variation in the extent to which they lend themselves to the verbal expression of aesthetic qualities.  Heterogeneity in the extent to which genres within domains (e.g. music or movies) lend themselves to verbal expression may explain why we observe educational stratification in terms of culture-mediated interaction within the domain in the absence of educational stratification in terms of overall levels of involvement.  For instance, as we saw above, members of all educational groups are about equally likely to report listening to music very often or fairly often; however, high education respondents are more likely to report {\em talking} about music with their friends.  
               
The extent to which the aesthetic qualities of a given cultural form is more or less likely to lend itself to verbal expression is of course not inherent in the form itself, but dependent on the institutional history of the cultural domain in question \citep{bourdieu93a}.  Drawing on recent theory and research on the phenomenon of {\em intellectulization} within artworlds \citep{baumann01,ferguson98}, we can hypothesize that the presence (or relative absence) of an already clearly specified, highly institutionalized critical vocabulary for culture-based talk, may explain variation in the extent to which we observe culture-mediated interaction in that genre.  In short, the fact that high education people are more likely to listen to musical genres that have these qualities, may explain their differential ability to draw on music as a cultural resource in conservation, even though they are no more likely to spend more time listening to music than low education respondents.
               
\subsubsection{Music}
Figure~\ref{fig4a} shows some evidence consistent with this hypothesis.  The figure shows the predicted probability of reporting discussing music with friends either ``fairly'' or ``very'' often, for persons who say that their favorite music is either ``rock and pop'' and ``country or folk,'' ``Christian or gospel,'' ``jazz and blues'' or ``classical or opera.'' These are coded as binary variables that equals one if the respondent says that one of these aggregate genre categories is among their favorite types of music.  The variables are not mutually exclusive, as each respondent is allowed to choose up to seven favorite genres\footnote{Most respondents have only one favorite. Only a small minority chooses more than two.}. The predicted probabilities comes for a logistic regression model that controls for education and all of the socio-demographic variables mentioned above.  Thus, they give us the adjusted effect of {\em taste} for particular genres on the probability of using music as a cultural resource in interaction.
               
Although the fact that reports on taste are forced to conform to this highly aggregate (and somewhat heterogeneous) genre categories is not ideal, the overall pattern of results conforms to expectations.  The musical genres most likely to have already institutionalized organizational and critical resources for intellectualization, are also the ones most likely to produce culture-based talk within the music domain. The genres most associated with either ``scenes'' or ``industry'' \citep{lena_peterson08} on the other hand are no more likely to generate conversation than expected by chance after controlling for the respondent's socio-demographic profile.  These results help to partly explain why we observe stratification of culture-mediated interaction in the music domain in the absence of education-linked consumption differences.

\subsubsection{Literature}          
The fact that the 2004 wave of National Social Survey included information on respondents favorite types of books allow us to further investigate the validity and the generality of the mechanism offered above to explain variation in the extent to which music becomes a cultural resource for conversation across educational groups. If it is correct that those subgenres for which there exist already established critical repertoires of discourse and schemes of evaluation which are more likely to come to generate culture-mediated interaction within that domain, then we should expect that in the literary domain that respondents who prefer genres associated with {\em fiction and ``literature''} in the traditional sense in which this is defined in the scholastic realm (e.g. novels) should be more likely to use books as a cultural resource to generate conversation with friends than respondent primarily attracted to non-fiction.  Within literary genres associated with fiction, we should expect less legitimate genres to be less likely to be used as a resource for conversational interaction than more legitimate genres (with legitimacy being defined for instance in terms of whether they get reviewed by critical gatekeepers in the mass media \citep{hirsch72}).
              
Figure~\ref{fig4b} shows evidence consistent with these expectations.  This figure is analogous to figure~\ref{fig4a}, except that I use respondents  taste for different book genres as a predictor of culture-mediated interaction within the literary realm.  The result show that three out of the four literary fiction genres are associated with a higher than average probability of talking to friends either ``fairly'' or ``very'' often about books.  The only fiction genre that can clear be seen to not lead to this outcome is the least legitimate and most critically debased genre: romance.  None of the genres associated with non-fiction can be considered (at least statistically) to be likely to generate culture-mediated interaction although both self-help books and books about history and politics come pretty close ($p=0.09$).  The case of self-help is itself interesting. As Illouz \citeyearpar{illouz03} has recently argued the self-help genre has become one of the key genres in which persons verbalize a large range of personal and social issues in what she refers to as a largely informal postmodern ``public sphere,'' and thus should be the non-fiction genre most likely to generate culture, which is precisely what we observe ($p=0.09$).     
      
\section{Discussion}
\subsection{Summary of the Results}
The goals of these paper have been twofold: first to highlight the theoretical and empirical importance of the phenomenon that I have referred to as {\em culture-mediated interaction}. Beginning with the Simmelian legacy on the importance of seemingly ``trivial'' conversational content for the emergence of ``sociability'' (or the ``play form'' of association) to the contemporary stress on ``entertainment talk'' as the cement that holds status groups together as micro-relational communities of acquaintance \citep{collins75, dimaggio87}, I have argued that the deployment of cultural knowledge in concrete interaction is one of the central, but least understood, processes in the field of culture stratification \citep{purcell07}.  In spite of its centrality, it has been a relatively under-studied phenomenon. This paper has attempted to redress this issue, but examining the socio-demographic distribution and cultural correlates of the tendency to talk to friends about culture consumption experiences.  
                
The results show that culture-mediated interaction is indeed a pervasive phenomenon, with a non-negligible majority of Americans shown to engage in some form of culture-mediated interaction using material from the cultural domains considered here.  However, in spite of its widespread distribution, culture-mediated interaction is a {\em stratified} phenomenon.  As expected, rates of culture-mediated interaction are higher among active culture consumers.  Because the educated and the young are more likely to more heavily engage in aesthetic consumption in the first place, I find that the prevalence of culture-mediated interaction is larger among the educated.  Interestingly enough, in realms associated with the popular arts, culture-mediated interaction appears to be stratified by education even in the absence of stratification by education in the rates of engagement with products from that cultural domain.  
                
This leads the finding that while the empirical association between education and culture-mediated interaction is readily explained by the larger probability of the educated to engage the more legitimate forms of cultural practice (e.g. the arts, reading), the association education and culture-mediated interaction in more popular domain cannot be so easily explained. Secondary analysis suggest that this may be due to heterogeneity in the specific genres that the educated consume within broad cultural domains.  In essence, the higher educated tend to be attracted genre groups that have already been institutionalized by the relevant cultural authorities as ``something to talk about'' and which have develop a conceptual apparatus of ``intellectualization'' and a critical vocabulary with which aesthetic experiences can be verbalized and thus become part of informal talk.
                
\subsection{Limitations of the analysis}
Naturally, the conclusions outlined above should be taken as preliminary suggestions and guidelines for future research rather than as firmly established propositions.  Obviously the most clear limitation of this study concerns the limited range of cultural domain included in the NSS for which information on the rate of culture-mediated interaction is available.  In particular, as befits an arts participation survey, questions about culture-mediated interaction are weighted toward the most legitimate forms of aesthetic practice but are silent when it comes to more popular culture consumption experiences \citep{dimaggio_mukhtar04}. 
                
In particular, past research suggests that culture-mediated interaction centered around sports \citep{erickson96}; Fashion and Cuisine \citep{dimaggio87} and in particular television \citep{fiske87} are a consistent generator of culture-based talk in formal  and informal settings (e.g. the workplace).  Thus, we can be certain that the rates of culture based talk reported above underestimate by a substantial amount its prevalence in the American population.  
                
The absence of information on culture-based talk in more popular domains may also lead to an overstatement of its connection to markers of stratification, such as educational credentials; although it is reasonable to suppose that just like in the case of music and movies, we may be likely to observe education-based stratification in rates of culture-mediated interaction in the absence of education-based stratification in consumption rates.  
                
It is also likely that including a wider range of cultural practices and domains may heavily qualify the finding that it is those more legitimate cultural pursuits within domains that generate the most culture-based talk.  It is likely for instance that within the realm of television-programming, it is those genres endowed with the least amount of critical legitimacy (e.g. reality television) that generate the most culture-based talk, especially among the young.   The very limited research done on ``sports talk'' \citep{purcell07} suggests however, that ``intellectualization'' in the realm of sports (e.g. produced by adopting the critical vocabulary produced by sports journalists) does facilitate sports-based talk among the educated, as suggested by the above analysis. 
                
The other major limitation of the present study is its relative glossing of the {\em context} of consumption as a mediator of the extent to which given aesthetic consumption experiences will generate culture-based talk. Some genres have built-in rituals of audience reception in which ``interaction'' is expected to occur {\em within} the consumption site and {\em during} the performance (e.g. a rock concert) and which thus carry little expectations that the cultural-knowledge gained therein will be transposed to other informal contexts as a generator of conversation.  Culture that is consumed according to socially constructed standards of audience ``propriety'' (e.g. sitting quietly in rapt attention) may in turn generate a lot of culture-based talk precisely because direct interaction among audience members (verbal and non-verbal) is frowned upon during the process of reception.  Cultural domains that exhibit wide variation in the context (and thus the expected audience conventions) of reception, may be seen to have a ``bi-modal'' distribution of their likelihood of being used for culture-based talk (like music in these data), with non-negligible proportions of active consumers reporting not using (as well as using) culture for that domain as a conversational resource.                
                
                   
\subsubsection{Why cultural matters may matter more than ``important'' matters}
The pervasiveness and social distribution of the phenomenon of culture-mediated interaction has important implications for the ways in which data about social connectivity are collected and conceptualized.  The two large-scale efforts to collect data about patterns of network connectivity in the American population have relied on the ``important matters'' name generator methodological tool \citep{burt84, marsden87, mcpherson_etal06}. The presumption here is that there is some sort of homology between the alleged {\em importance of the content of conversation and the importance of the tie in question}.  Importance here being conceptualized both as 1) theoretical importance for the social scientist,--e.g. the status of the social ties identified in this manner as an measure of the extent to which the individual is ``connected'' \citep{burt90}; and 2) practical importance for the person; or the extent to which the connection identified in this manner give us access to a meaningful subset of the relations that a person would draw on in time of emotional or financial help \citep{wellman_wortley90}.
                  
But given the social distribution of culture -mediated interaction---and its intuitive connection with certain socio-demographic markers---it is perilous to suppose that the ``importance'' of the ``matters'' discussed has any connection to the ``importance'' of the contact with which they are discussed.  In fact, given the theoretical framework outlined here, we should expect little or no connection between the objective importance of the network content and the object (or social-scientific) importance of the ties connected to to those contents.  If Simmelian ``pure sociability'' occurs within those relationships which provide {\em the most intrinsic satisfaction} we should expect the connection to even be negative (e.g. meaningful and subjectively important ties are sustained through seemingly ``trivial'' and ``unimportant'' cultural content).  This becomes more important when we realize that, as Simmel would have predicted, conversation involving ``important matters'' is too sporadic and inconsistent a phenomenon with which to judge whether persons are connected or not \citep{bearman_parigi04}.
                  
A substantial number of people (including a larger proportion of married persons) report simply not having anybody with which to discuss ``important'' matters with \citep{bearman_parigi04}). Should we conclude therefore that most persons are ``disconnected'' and experience little sociability \citep{mcpherson_etal06, putnam00}?  I would suggest that a more theoretically sophisticated understanding of what ``sociability'' is all about suggests otherwise. ``Important matters'' name generators may not be getting at either ``important'' content  or important ties (if the social-scientific interest is to uncover those ties that are meaningful for the individual).  This becomes even more of an issue when we consider that interaction based on important matters, may be more steeply stratified by markers of social status, such as education, than culture-based talk \citep{bourdieu84, mcpherson_etal06, marsden87, fischer82}.
                  
In this respect, the ``important matters'' name-generator appears to be a tough ``test'', that only a few privileged persons endowed with certain forms of linguistic (or political) capital may pass \citep{bourdieu84}. Important matters name generators may thus confound stratification based on aptitudes orthogonal to social competence and ability to form connections with stratification criteria that directly impinge on the persons' ability to in fact create and sustain social relationships \citep{lin99, dimaggio96, lizardo06a}
                  
If the framework outlined here is on the right track, {\em cultural matters may matter more than important matters} as the content with which to sustain (and thus measure) meaningful ties.  Thus, querying persons as to whom they discuss the ``trivial'' matters related to the cultural content that makes up entertainment talk (e.g. sports, fashion, cuisine, television, movies, the arts), maybe a more sensible strategy to generate data on personal social connectivity, and to  draw conclusions about the fate of this connectivity in the contemporary situation.

\newpage
\bibliographystyle{asr}
\bibliography{N:/Private/zotero-library}

\pagebreak
\begin{landscape}
   \begin{table}[htb!]
           \begin{tabular}{l*{13}{c}}
           \hline \hline
           
\multicolumn{13}{c}{Culture-Mediated Interaction}\\
\hline
&\multicolumn{3}{c}{Education}&\multicolumn{3}{c}{Age}&\multicolumn{3}{c}{Gender}&\multicolumn{3}{c}{Model}\\
\hline
& $\chi^{2}$  & {\em p} & df & $\chi^{2}$ & {\em p} & df & $\chi^{2}$ &  {\em p} & df &  $\chi^{2}$ &  Educ \% &  Age \%  &  Gender \% \\ 
Movies&47.0&0.0000&4&202.0&0.0000&6&22.5&0.0000&1&312.8&15.0&64.6&7.2\\
Books&93.9&0.0000&4&13.2&0.0398&6&60.6&0.0000&1&171.6&54.7&7.7&35.3\\
Music&81.6&0.0000&4&56.9&0.0000&6&2.1&0.1450&1&157.5&51.8&36.1&1.3\\
Plays&105.8&0.0000&4&19.8&0.0030&6&24.0&0.0000&1&173.3&61.1&11.4&13.9\\
Arts&105.3&0.0000&4&12.4&0.0531&6&14.5&0.0001&1&146.4&71.9&8.5&9.9\\
\hline
\hline
\multicolumn{13}{c}{Aesthetic Consumption}\\
\hline
&\multicolumn{3}{c}{Education}&\multicolumn{3}{c}{Age}&\multicolumn{3}{c}{Gender}&\multicolumn{3}{c}{Model}\\
\hline
& Educ $\chi^{2}$  & {\em p} & df & Age $\chi^{2}$ & {\em p} & df & Gender $\chi^{2}$ &  {\em p} & df &  Model $\chi^{2}$ &  Educ \% &  Age \%  &  Gender \% \\ 
Movies&56.0&0.0000&4&175.7&0.0000&6&4.2&0.0403&1&327.9&17.1&53.6&1.3\\
Books&70.4&0.0000&4&10.3&0.1132&6&35.6&0.0000&1&123.9&56.8&8.3&28.7\\
Music&4.3&0.3663&4&77.1&0.0000&6&6.4&0.0113&1&146.3&2.9&52.7&4.4\\
Plays&152.2&0.0000&4&5.9&0.4319&6&3.2&0.0729&1&181.5&83.9&3.3&1.8\\
Arts&112.5&0.0000&4&19.2&0.0038&6&1.2&0.2809&1&169.3&66.4&11.3&0.7\\
            \hline
           \end{tabular}
           \caption{}
           \label{tab1}
   \end{table}
\end{landscape}


   \begin{table}[htb!]   									\begin{tabular*}{\hsize}{@{\hskip\tabcolsep\extracolsep\fill}l l*{3}{c}} 
		\hline \hline	     
	     \multicolumn{4}{c}{Movies (Age)}\\ 
	     \hline 
	     Model&$G^{2}$&{\em df}&BIC\\
	     Spurious Association by Social Position&484.87&63&187.60\\
		Interaction as Intervening&172.23&72&-167.50\\
		Taste as Intervening&173.72&72&-166.01\\
		No Three-Way Interaction&74.36&54&-180.44*\\
		\hline \hline
		\multicolumn{4}{c}{Music (Age)}\\ 
		\hline
		Model&$G^{2}$&{\em df}&BIC\\
		Spurious Association by Social Position&227.41&63&-69.86\\
		Interaction as Intervening&160.81&72&-178.92\\
		Taste as Intervening&98.72&72&-241.01*\\
		No Three-Way Interaction&60.99&54&-193.80\\
		\hline \hline
		\multicolumn{4}{c}{Books (Education)}\\ 
		\hline
		Model&$G^{2}$&{\em df}&BIC\\
		Spurious Association by Social Position&469.95&45&272.76\\
		Interaction as Intervening&72.58&48&-137.76*\\
		Taste as Intervening&83.40&48&-126.93\\
		No Three-Way Interaction&31.79&36&-125.97\\
		\hline \hline
		\multicolumn{4}{c}{Books (Gender)}\\ 
		\hline
		Model&$G^{2}$&{\em df}&BIC\\	
		Spurious Association by Social Position&464.36&18&401.98\\
		Interaction as Intervening&16.02&12&-25.57*\\
		Taste as Intervening&30.71&12&-10.88\\
		No Three-Way Interaction&5.26&9&-25.93\\
		\hline \hline
		\multicolumn{4}{c}{Arts (Education)}\\ 
		\hline
		Model&$G^{2}$&{\em df}&BIC\\
		Spurious Association by Social Position&582.56&45&385.37\\
		Interaction as Intervening&138.37&48&-71.96\\
		Taste as Intervening&79.80&48&-130.53*\\
		No Three-Way Interaction&49.76&36&-107.99\\
		\hline \hline
		\multicolumn{4}{c}{Plays (Education)}\\ 
		\hline
		Model&$G^{2}$&{\em df}&BIC\\
		Spurious Association by Social Position&443.40&45&246.21\\
		Interaction as Intervening&85.29&48&-125.05\\
		Taste as Intervening&68.19&48&-142.15*\\
		No Three-Way Interaction&26.68&36&-131.07\\

		\hline
		\multicolumn{3}{l}{*Preferred Model}\\ 
        \end{tabular*}
        \caption{}
        \label{tab2}
   \end{table}

\newpage
   \begin{figure}[here]
       \centering
       \subfigure[Univariate distribution of culture-mediated interaction probabilities (whole sample)]{
                \includegraphics[width=7 in, height=4 in]{univariate-frequency-talk}
                \label{fig1a}
                }
         \subfigure[Univariate distribution of culture-mediated interaction probabilities (culturally active subsample)]{         
                \includegraphics[width=7 in, height=4 in]{univariate-dist-hi-consump}
                \label{fig1b}
                }     
        \caption{Univariate distribution of culture mediated interaction, 2002 and 2004 National Social Survey, Arts and Culture Supplement}
  \end{figure}   

\newpage
   \begin{figure}[here]
       \centering
       \subfigure[Predicted probabilities of culture-mediated interaction, by education]{
                \includegraphics[width=7 in, height=4 in]{educ-pred-probs-interact-binary}
                \label{fig2a}
                }
         \subfigure[Predicted probabilities of aesthetic consumption, by education]{         
                \includegraphics[width=7 in, height=4 in]{educ-pred-probs-consump-binary}
                \label{fig2b}
                }     
        \caption{Predicted probabilities of discussing cultural matters with friends and engaging in culture consumption for five cultural domains, 2002 and 2004 National Social Survey, Arts and Culture Supplement}
  \end{figure}   
  
 \newpage
   \begin{figure}[here]
       \centering
       \includegraphics[width=6 in, height=8 in]{educ-age-music-effect}    
        \caption{Predicted probabilities of discussing music with friends very often by Education and Age, 2002 and 2004 National Social Survey, Arts and Culture Supplement}
        \label{fig3}
  \end{figure}  


   \begin{figure}[here]
      \centering
         \includegraphics[width=6.5 in, height=5 in]{mfx-educ-pred-probs}
         \caption{Marginal effects of educational transitions on the probabilities of engaging in frequent culture-mediated interaction before (reduced model) and after (full model) controlling for consumption frequencies.}
         \label{fig4}
   \end{figure}

\newpage
   \begin{figure}[here]
       \centering
       \subfigure{
                \includegraphics[width=5 in, height=3 in]{music-interaction-by-genre}
                \label{fig5a}
                }
         \subfigure{         
                \includegraphics[width=5 in, height=3 in]{book-interaction-by-genre}
                \label{fig5b}
                }     
        \caption{}
  \end{figure}   
\end{document}
